\VerdierPhysicalPageAnchor{0063}%
\VerdierSemanticAnchor{ch1-1-6-5-continuation}
d'où on déduit immédiatement :

\VerdierSemanticAnchor{ch1-1-6-5-deduction}
\VerdierSemanticAnchor{ch1-1-6-5-beta-s-coboundary}
\begin{equation*}
  \alpha_s(\sigma_1,\sigma_2)\alpha(\sigma_1,\sigma_2)^{-1}
  =\beta_s(\sigma_2)\beta_s(\sigma_1+\sigma_2)^{-1}\beta_s(\sigma_1),
\end{equation*}

\VerdierSemanticAnchor{ch1-1-6-5-eq3}
\VerdierSemanticAnchor{ch1-1-6-5-eq3-beta}
\VerdierSemanticAnchor{ch1-1-6-5-eq3-basis}
\begin{equation}
  \begin{aligned}
    \beta_s(\sigma_1)
      &=(-1)^{\left(
          \sum_{\substack{i<j\\s(i)>s(j)}}p_i p_j
        \right)},\\
    \beta_s(e_i)&=1,
      \qquad \forall i\in[1,n].
  \end{aligned}
  \tag{1.6.5.3}
\end{equation}

\VerdierSemanticAnchor{ch1-1-6-5-application}
Par suite, d'après (1.6.2), toute application
$\beta\colon\mathbb Z^n\longrightarrow\mathbb Z/2\mathbb Z$ telle que

\VerdierSemanticAnchor{ch1-1-6-5-eq4}
\begin{equation}
  \alpha_s(\sigma_1,\sigma_2)\alpha(\sigma_1,\sigma_2)^{-1}
  =\beta(\sigma_2)\beta(\sigma_1+\sigma_2)^{-1}\beta(\sigma_1)
  \tag{1.6.5.4}
\end{equation}

\VerdierSemanticAnchor{ch1-1-6-5-eq5-intro}
est du type :

\VerdierSemanticAnchor{ch1-1-6-5-eq5}
\begin{equation}
  \beta=\beta_s\eta,
  \tag{1.6.5.5}
\end{equation}

\VerdierSemanticAnchor{ch1-1-6-5-closure}
où $\eta\colon\mathbb Z^n\longrightarrow\mathbb Z/2\mathbb Z$ est
l'unique homomorphisme de groupes qui prend sur les
$e_j\in\mathbb Z^n$ les valeurs $\beta(e_j)$.

\VerdierSemanticAnchor{ch1-1-6-6}
\VerdierSemanticAnchor{ch1-1-6-6-product-setup}
\noindent\textbf{1.6.6.}\quad
Soient $(C_i,\mathbb Z,T_i,c_i)$, $1\leq i\leq n$, des
$\mathbb Z$-catégories et $(D,\mathbb Z,T,c)$ une $\mathbb Z$-catégorie
additive. Désignons par $(C,\mathbb Z^n,T,c)$ la catégorie produit des
$(C_i,\mathbb Z,T_i,c_i)$. Le groupe $\mathbb Z^n$ est muni
naturellement d'une base ordonnée $e_1,\ldots,e_n$. Il opère à
isomorphisme près sur $C$ par la formule (on utilise ici la notation à
droite introduite en (1.3.12)) :

\VerdierSemanticAnchor{ch1-1-6-6-eq1}
\begin{equation}
  (X_1,\ldots,X_n)[\sigma]
  =(X_1[l_1],\ldots,X_n[l_n]),
  \qquad
  \sigma=\sum_{i=1}^{n}l_i e_i.
  \tag{1.6.6.1}
\end{equation}

\VerdierSemanticAnchor{ch1-1-6-6-phi-action}
Soit $\varphi\colon\mathbb Z^n\longrightarrow\mathbb Z$
l'homomorphisme
$\varphi(\sum_{i=1}^{n}l_i e_i)=\sum_{i=1}^{n}l_i$.
Faisons opérer $\mathbb Z^n$ sur $D$ par l'intermédiaire de $\varphi$
(1.4.1). Dans cette situation, nous utiliserons toujours les
conventions suivantes :

\VerdierSemanticAnchor{ch1-1-6-6-sign-rule-1}
\noindent\textbf{Règle des signes I :}
\emph{Les foncteurs de $C$ dans $D$ seront toujours, sauf mention
expresse du contraire, des $\mathbb Z^n$-foncteurs tordus par le
cocycle de Koszul relatif à la base ordonnée $e_1,\ldots,e_n$.}

\VerdierSemanticAnchor{ch1-1-6-6-sign-rule-2}
\noindent\textbf{Règle des signes II :}
\emph{les morphismes de foncteurs seront, sauf mention expresse du
contraire, des morphismes de $\mathbb Z^n$-foncteurs relatifs à
l'homomorphisme trivial
$\beta\colon\mathbb Z^n\longrightarrow\mathbb Z/2\mathbb Z$ (1.4.6).}
